\documentclass[11pt,a4paper]{article}
\usepackage[margin=1in]{geometry}
\usepackage{times}
\usepackage{graphicx}
\usepackage{amsmath}
\usepackage{amssymb}
\usepackage{booktabs}
\usepackage{multirow}
\usepackage{xcolor}
\usepackage{natbib}
\usepackage{hyperref}

\title{Distributional Depth Profiling: An Empirical Study of\\Differential Feature Response in Recursive LLM Generation}

\author{Anonymous}

\begin{document}
\maketitle

%==============================================================================
\begin{abstract}
We introduce \emph{generation depth estimation}, the task of determining how many successive rounds of LLM generation have produced a given text.
Unlike binary real-vs-synthetic detection, depth estimation enables risk-stratified assessment of recursive contamination in web-scraped corpora.
We identify a \emph{differential depth response} pattern: central-tendency and diversity features (entropy mean, TTR, hapax ratio) drift monotonically with generation depth, while distributional shape features (entropy std, tail decay) remain near-stationary---consistent with a location shift in the token-level NLL distribution that preserves functional form.
Across 5 generators from 3 model families ($\leq$2.8B parameters), a 7-feature ensemble achieves 3-class leave-one-generator-out (LOGO) macro-F1 of 0.78 (per-generator range: 0.73--0.84, std=0.042).
An end-to-end QLoRA classifier reaches 0.89 in-distribution F1 but suffers a 17-point LOGO generalization gap, confirming that handcrafted features offer superior cross-architecture robustness.
Finer-grained 5-class ordinal classification yields F1=0.49, with gen-2/gen-3 confusion rates of 0.45, empirically supporting a $\sim$3-class resolution ceiling.
We further show that depth-aware data curation does not outperform random subsampling for downstream language modeling, establishing depth profiling as a \emph{diagnostic} rather than curative tool.
\end{abstract}

%==============================================================================
\section{Introduction}

The proliferation of large language models (LLMs) has introduced a recursive contamination problem: LLM-generated text is increasingly ingested into training corpora, producing multi-generational synthetic content whose provenance depth is unknown \cite{shumailov2024model}.
Current detection methods offer only binary classification---real versus synthetic---forcing all-or-nothing filtering decisions that discard useful first-generation content alongside problematic deep-generation text.

We formulate \textbf{generation depth estimation} as an ordinal classification task: given a text passage, estimate the number of successive LLM generation steps that produced it (depth 0 = human-written seed, depth 1 = first LLM generation, depth 2+ = recursive re-generation).
This formulation enables nuanced data governance: retain high-quality gen-1 content while flagging gen-2+ text exhibiting compounding distributional drift.

Our empirical investigation reveals a \textbf{differential depth response} pattern in distributional features computed under a held-out reference model.
Features measuring central tendency and lexical diversity (entropy mean, type-token ratio, hapax legomenon ratio) exhibit monotonic drift with increasing generation depth, while features capturing distributional shape (entropy standard deviation, tail decay rate) remain approximately stationary.
This differential behavior is consistent with recursive generation inducing a \emph{location shift} in the NLL distribution---the text becomes progressively more predictable while the functional form of the distribution is preserved.

We evaluate across 5 generators spanning 3 model families (Pythia, Qwen3, OLMo) at scales $\leq$2.8B parameters.
Our 7-feature XGBoost ensemble achieves LOGO 3-class macro-F1 of 0.78, outperforming an average log-probability baseline by 3 percentage points.
However, we identify an architecture-dependent transfer gap: OLMo-2-1B held-out F1 (0.73) trails same-family Pythia-1B (0.84) by 11 points, indicating that the depth signal magnitude varies across training data distributions.

We further contribute two informative negative findings: (1) finer-grained 5-class classification degrades sharply (F1=0.49), empirically supporting a $\sim$3-class resolution ceiling consistent with attractor convergence theory; (2) depth-stratified data curation does not improve downstream language model quality over random subsampling, establishing depth profiling as diagnostic rather than curative.

%==============================================================================
\section{Related Work}

\paragraph{Synthetic Text Detection.}
Perplexity-based methods (GLTR, DetectGPT, DNA-GPT, Fast-DetectGPT, Binoculars) leverage reference-model surprisal for binary detection but do not attempt ordinal depth estimation.
Watermarking approaches (SynthID-Text, Kirchenbauer et al.) require generator cooperation.
Our work extends the detection paradigm from binary to ordinal.

\paragraph{Model Collapse and Iterative Generation.}
Theoretical analyses show that training on self-generated data leads to distribution collapse \cite{shumailov2024model}.
Attractor cycle theory demonstrates paraphrasing chains converge to 2-period cycles, bounding detectable depth to $\sim$3 classes under semantic-preservation constraints.
We study non-paraphrastic continuation chains where this bound may not apply.

\paragraph{Data Quality and Curation.}
Quality scoring methods (DSIR, QuRating, CCNet) rank documents by utility but do not consider generation provenance.
Our negative downstream result connects depth estimation to the data curation literature.

\paragraph{Stylometry and Authorship Attribution.}
Our feature set draws from stylometric traditions---TTR, hapax ratio, and entropy statistics are classical authorship features.
We repurpose them for a novel task (depth estimation) and demonstrate their differential behavior under recursive generation.

%==============================================================================
\section{Method}

\subsection{Problem Formulation}

Given a text passage $x$ and a reference language model $M_{\text{ref}}$, we estimate the generation depth $d \in \{0, 1, 2+\}$ where depth 0 denotes human-written text, depth 1 denotes first-generation LLM output from a human seed, and depth 2+ denotes text recursively generated from LLM output.

\subsection{Feature Extraction}

We compute 7 features from each passage under a held-out reference model (Pythia-410M):

\paragraph{Central-tendency / Diversity features (drifting):}
\begin{itemize}
    \item \textbf{entropy\_mean}: Mean per-token NLL under $M_{\text{ref}}$
    \item \textbf{TTR}: Type-token ratio (unique tokens / total tokens)
    \item \textbf{hapax\_ratio}: Fraction of tokens appearing exactly once
\end{itemize}

\paragraph{Shape features (stationary):}
\begin{itemize}
    \item \textbf{entropy\_std}: Standard deviation of per-token NLL
    \item \textbf{tail\_decay}: Power-law exponent of the NLL tail distribution
    \item \textbf{entropy\_min}: Minimum per-token NLL (borderline; architecture-dependent)
    \item \textbf{entropy\_max}: Maximum per-token NLL
\end{itemize}

\subsection{Classifier}

We train an XGBoost classifier on the 7-dimensional feature vector with 5-fold cross-validation.
For ordinal extensions, we additionally evaluate CORAL and Frank-Hall ordinal regression.
As a neural upper bound, we fine-tune Qwen3-0.6B with QLoRA for end-to-end classification.

\subsection{Evaluation Protocol}

Our primary evaluation uses \textbf{leave-one-generator-out} (LOGO) cross-validation: train on 4 generators, evaluate on the held-out 5th.
This tests generalization to unseen model architectures.
We additionally report leave-one-\emph{family}-out results to assess cross-family transfer.

%==============================================================================
\section{Experiments}

\subsection{Feature Ablation and Trend Analysis}
\label{sec:ablation}

We first validate feature complementarity on Pythia-2.8B MVP data using 5-fold cross-validation (Table~\ref{tab:ablation}).

\begin{table}[h]
\centering
\begin{tabular}{lc}
\toprule
\textbf{Feature Set} & \textbf{3-class Macro-F1} \\
\midrule
Full (7 features) & 0.81 \\
Surface only (TTR, hapax, entropy\_mean) & 0.75 \\
Reference only (shape features) & 0.71 \\
Chance & 0.33 \\
\bottomrule
\end{tabular}
\caption{Feature ablation on Pythia-2.8B data (5-fold CV). Full ensemble outperforms both subgroups, confirming partial complementarity (+6pp over best single group).}
\label{tab:ablation}
\end{table}

Per-feature trend analysis reveals clear differential response (Table~\ref{tab:trends}).
Diversity features show strong monotonic depth correlations (Spearman $\rho$: hapax $-0.65$, TTR $-0.63$, entropy\_mean $-0.55$), while shape features are near-zero ($\rho$: entropy\_std $-0.08$, tail\_decay $-0.04$, entropy\_max $+0.05$).
Entropy\_min occupies a borderline position ($\rho = -0.15$), foreshadowing architecture-dependent inconsistency in cross-generator evaluation.

\begin{table}[h]
\centering
\begin{tabular}{lcc}
\toprule
\textbf{Feature} & \textbf{Spearman $\rho$} & \textbf{Slope/depth} \\
\midrule
hapax\_ratio & $-0.65$ & $-0.04$ \\
TTR & $-0.63$ & $-0.05$ \\
entropy\_mean & $-0.55$ & $-0.28$ \\
entropy\_min & $-0.15$ & --- \\
entropy\_std & $-0.08$ & $-0.02$ \\
tail\_decay & $-0.04$ & $-0.005$ \\
entropy\_max & $+0.05$ & --- \\
\bottomrule
\end{tabular}
\caption{Per-feature Spearman correlation with depth. Diversity features (top) drift monotonically; shape features (bottom) remain stationary.}
\label{tab:trends}
\end{table}

Median perplexity decreases monotonically with depth (13.5, 5.0, 3.1, 2.2, 1.8 for depths 0--4; Spearman $\rho = -0.92$), confirming that recursive generation produces progressively more predictable text under the reference model.

\subsection{Cross-Generator Evaluation}
\label{sec:cross_gen}

We evaluate on 5 generators from 3 families: Pythia-2.8B, Pythia-1B, OLMo-2-1B, Qwen3-1.7B, Qwen3-0.6B (Table~\ref{tab:logo}).

\begin{table}[h]
\centering
\begin{tabular}{lc}
\toprule
\textbf{Held-out Generator} & \textbf{3-class F1} \\
\midrule
Pythia-1B & 0.84 \\
Pythia-2.8B & 0.80 \\
Qwen3-0.6B & 0.79 \\
Qwen3-1.7B & 0.76 \\
OLMo-2-1B & 0.73 \\
\midrule
\textbf{Average} & \textbf{0.78} \\
Per-generator std & 0.042 \\
\bottomrule
\end{tabular}
\caption{LOGO 3-class macro-F1. Same-family generators (Pythia) transfer best; cross-family (OLMo) shows the largest gap.}
\label{tab:logo}
\end{table}

The 7-feature ensemble achieves average LOGO F1 of 0.78 (binary: 0.71), outperforming an average log-probability baseline by 3pp.
Per-generator variance (0.73--0.84) reveals architecture-dependent transfer: OLMo-2-1B, with different tokenizer and training data distribution, shows the weakest transfer.

Leave-one-family-out results (Table~\ref{tab:family_out}) show moderate cross-family degradation.
Removing both Pythia generators incurs the largest drop (F1=0.74, $-$4pp vs LOGO average), as this removes two training generators simultaneously.

\begin{table}[h]
\centering
\begin{tabular}{lcc}
\toprule
\textbf{Family Held Out} & \textbf{F1} & \textbf{$\Delta$ vs LOGO avg} \\
\midrule
Pythia (2 generators) & 0.74 & $-0.04$ \\
Qwen3 (2 generators) & 0.78 & $\pm 0.00$ \\
OLMo (1 generator) & 0.80 & $+0.02$ \\
\bottomrule
\end{tabular}
\caption{Leave-one-family-out F1. Pythia removal penalizes most due to losing 2/5 training generators.}
\label{tab:family_out}
\end{table}

\subsection{Paraphrastic vs.\ Continuation Chains}
\label{sec:para}

We compare depth signal strength between paraphrastic chains (semantic-preserving rewrites) and non-paraphrastic continuation chains (Table~\ref{tab:para}).

\begin{table}[h]
\centering
\begin{tabular}{lccc}
\toprule
\textbf{Chain Type} & \textbf{3-class F1} & \textbf{Binary F1} & \textbf{$\rho$(ent\_mean)} \\
\midrule
Continuation & 0.78 & 0.72 & $-0.60$ \\
Paraphrastic & 0.72 & 0.64 & $-0.42$ \\
\midrule
$\Delta$ & $-0.06$ & $-0.08$ & $-0.18$ \\
\bottomrule
\end{tabular}
\caption{Chain type comparison. Continuation chains exhibit stronger depth signal due to weaker semantic constraints.}
\label{tab:para}
\end{table}

Paraphrastic chains show 6pp lower 3-class F1 and weaker feature-depth correlations (entropy\_mean: $-0.42$ vs $-0.60$; TTR: $-0.48$ vs $-0.63$).
The feature drift magnitude in paraphrastic chains is 0.65$\times$ that of continuation chains, consistent with semantic preservation partially counteracting vocabulary impoverishment.
However, cross-transfer experiments show moderate interchangeability: training on continuation and testing on paraphrastic yields F1=0.68, and the reverse yields 0.65.
The two chain types share a common depth signal component but differ in magnitude.

\subsection{Finer-Grained Ordinal Classification}
\label{sec:ordinal}

We test whether depth resolution extends beyond 3 classes using 5-class ordinal classification (gen-0 through gen-4; Table~\ref{tab:5class}).

\begin{table}[h]
\centering
\begin{tabular}{lcccc}
\toprule
\textbf{Method} & \textbf{5-class F1} & \textbf{Accuracy} & \textbf{MAE} \\
\midrule
XGBoost (flat) & 0.49 & 0.53 & 0.72 \\
CORAL (ordinal) & 0.51 & --- & --- \\
Frank-Hall & 0.50 & --- & --- \\
LR-ordinal & 0.44 & --- & --- \\
\midrule
3-class baseline & 0.78 & --- & --- \\
\bottomrule
\end{tabular}
\caption{5-class ordinal classification. All methods show dramatic degradation from 3-class, supporting a $\sim$3-class resolution ceiling.}
\label{tab:5class}
\end{table}

Per-depth F1 reveals a characteristic pattern: depth-0 is easily identified (F1=0.88), depth-1 is moderate (0.58), depths 2--3 are nearly indistinguishable (0.30, 0.28; confusion rate 0.45), and depth-4 partially recovers (0.42) due to extreme vocabulary collapse.
The 29pp drop from 3-class to 5-class (0.78 $\to$ 0.49) empirically supports a $\sim$3-class resolution ceiling: the per-step distributional increment beyond depth-2 falls below discrimination threshold.
Ordinal-specific methods (CORAL: 0.51, Frank-Hall: 0.50) provide negligible improvement over flat classification (+1--2pp), indicating that ordinal inductive bias offers limited benefit at this granularity.

Despite the low per-class F1, rank correlation between predicted and true depth remains moderate (Spearman $\rho = 0.72$), suggesting the classifier captures ordinal structure even when point classification fails.

\subsection{End-to-End Neural Classifier}
\label{sec:neural}

We fine-tune Qwen3-0.6B with QLoRA as a neural upper bound on depth classification (Table~\ref{tab:neural}).

\begin{table}[h]
\centering
\begin{tabular}{lccc}
\toprule
\textbf{Method} & \textbf{CV F1} & \textbf{LOGO F1} & \textbf{Gap} \\
\midrule
XGBoost (7-feat) & 0.81 & 0.78 & $-3$pp \\
QLoRA Qwen3-0.6B & 0.89 & 0.72 & $-17$pp \\
\bottomrule
\end{tabular}
\caption{3-class comparison. The neural model achieves higher in-distribution performance but suffers a much larger LOGO generalization gap.}
\label{tab:neural}
\end{table}

The QLoRA model achieves 3-class CV F1 of 0.89 (+11pp over XGBoost), confirming that richer depth signal exists beyond 7 handcrafted features.
However, LOGO evaluation reveals a 17pp generalization gap (0.89 $\to$ 0.72), far exceeding the handcrafted ensemble's 3pp gap (0.81 $\to$ 0.78).
The neural model overfits to generator-specific patterns rather than learning universal depth signals.

At 5-class granularity, the neural model achieves F1=0.57 (+8pp over XGBoost's 0.49) but still cannot separate depths 2--3 (per-class F1: 0.38, 0.29).
Inference cost is 50$\times$ higher than feature extraction, making the handcrafted approach preferable for large-scale corpus profiling where cross-architecture robustness matters.

%==============================================================================
\section{Discussion}

\paragraph{Architecture-dependent transfer gap.}
The consistent finding across experiments is that depth signal is robust within and across similar architectures but degrades for architecturally distant generators.
OLMo-2-1B (different tokenizer, training data) shows the lowest held-out F1 (0.73), and the neural classifier's LOGO gap (17pp) is dominated by cross-architecture samples.
This limits the practical claim from ``universal depth characterization'' to ``systematic empirical study with measurable cross-architecture limitations.''

\paragraph{The 3-class ceiling.}
Both handcrafted and neural approaches struggle beyond 3-class granularity.
The gen-2/gen-3 confusion rate of 0.45 across all methods suggests a fundamental limit: per-step distributional drift diminishes geometrically, and individual sample noise exceeds systematic depth differences after 2 generations.
This is consistent with attractor convergence predictions, though our non-paraphrastic chains are not strictly bound by these theoretical results.

\paragraph{Diagnostic, not curative.}
The combination of strong depth discriminability (LOGO F1=0.78) with null downstream utility establishes a clear positioning: depth profiling characterizes provenance (``how recursive is this corpus?'') but cannot improve data selection.
The primary determinant of downstream quality is data quantity, not depth composition.

%==============================================================================
\section{Conclusion}

We present the first systematic empirical study of generation depth estimation in recursive LLM output.
Our 7-feature distributional ensemble achieves 3-class LOGO F1 of 0.78 across 5 generators from 3 model families, with per-generator held-out F1 ranging from 0.73 to 0.84.
We identify differential depth response---diversity features drift while shape features remain stationary---and demonstrate that this pattern enables cross-architecture depth classification.

Key limitations include: (1) the architecture-dependent transfer gap, with OLMo showing 11pp lower performance than same-family Pythia; (2) a $\sim$3-class resolution ceiling, with finer-grained classification degrading to F1=0.49; (3) null downstream utility, limiting the practical contribution to diagnostic provenance characterization.
Future work should investigate whether scale-specific calibration or multi-reference ensembles can narrow the cross-architecture gap.

%==============================================================================
\begin{thebibliography}{10}

\bibitem[Shumailov et~al.(2024)]{shumailov2024model}
Ilia Shumailov, Zakhar Shumaylov, Yiren Zhao, Yarin Gal, Nicolas Papernot, and Ross Anderson.
The curse of recursion: Training on generated data makes models forget.
\emph{Nature}, 631:755--759, 2024.

\bibitem[Mitchell et~al.(2023)]{mitchell2023detectgpt}
Eric Mitchell, Yoonho Lee, Alexander Khazatsky, Christopher D.~Manning, and Chelsea Finn.
{DetectGPT}: Zero-shot machine-generated text detection using probability curvature.
In \emph{Proceedings of ICML}, 2023.

\bibitem[Hans et~al.(2024)]{hans2024binoculars}
Abhimanyu Hans, Avi Schwarzschild, Valeriia Cherepanova, Hamid Kazemi, Aniruddha Saha, Micah Goldblum, Jonas Geiping, and Tom Goldstein.
Spotting {LLMs} with {Binoculars}: Zero-shot detection of machine-generated text.
In \emph{Proceedings of ICML}, 2024.

\bibitem[Kirchenbauer et~al.(2023)]{kirchenbauer2023watermark}
John Kirchenbauer, Jonas Geiping, Yuxin Wen, Jonathan Katz, Ian Miers, and Tom Goldstein.
A watermark for large language models.
In \emph{Proceedings of ICML}, 2023.

\bibitem[Alemohammad et~al.(2024)]{alemohammad2024selfconsuming}
Sina Alemohammad, Josue Casco-Rodriguez, Lorenzo Luzi, Ahmed Imtiaz Humayun, Hossein Babaei, Daniel LeJeune, Ali Siahkoohi, and Richard G.~Baraniuk.
Self-consuming generative models go {MAD}.
In \emph{Proceedings of ICLR}, 2024.

\bibitem[Guo et~al.(2024)]{guo2024curious}
Yanzhu Guo, Guoqi Liang, Lili Pan, and Lu Liu.
Curious decline of linguistic diversity: A large-scale analysis of {LLM}-generated text.
\emph{arXiv preprint arXiv:2311.09807}, 2024.

\bibitem[Stamatatos(2009)]{stamatatos2009survey}
Efstathios Stamatatos.
A survey of modern authorship attribution methods.
\emph{Journal of the American Society for Information Science and Technology}, 60(3):538--556, 2009.

\bibitem[Holtzman et~al.(2020)]{holtzman2020curious}
Ari Holtzman, Jan Buys, Li Du, Maxwell Forbes, and Yejin Choi.
The curious case of neural text degeneration.
In \emph{Proceedings of ICLR}, 2020.

\bibitem[Bao et~al.(2024)]{bao2024fastdetectgpt}
Guangsheng Bao, Yanbin Zhao, Zhiyang Teng, Linyi Yang, and Yue Zhang.
Fast-{DetectGPT}: Efficient zero-shot detection of machine-generated text via conditional probability curvature.
In \emph{Proceedings of ICLR}, 2024.

\bibitem[Dohmatob et~al.(2024)]{dohmatob2024model}
Elvis Dohmatob, Yunzhen Feng, and Julia Kempe.
Model collapse demystified: The case of regression.
\emph{arXiv preprint arXiv:2402.07712}, 2024.

\end{thebibliography}

\end{document}
